\documentclass[10pt,a4paper]{amsart}
\usepackage{amsmath,amssymb,amsthm,graphicx}
\usepackage{verbatim}
\newenvironment{reference}{\comment}{\endcomment}
%\newenvironment{reference}{}{}
\newenvironment{slogan}{\comment}{\endcomment}
\newenvironment{history}{\comment}{\endcomment}
% For commutative diagrams we use Xy-pic
\usepackage[all]{xy}

% We use 2cell for 2-commutative diagrams.
\xyoption{2cell}
% Theorem environments.
%
\theoremstyle{plain}
\newtheorem{theorem}[subsection]{Theorem}
\newtheorem{proposition}[subsection]{Proposition}
\newtheorem{lemma}[subsection]{Lemma}

\theoremstyle{definition}
\newtheorem{definition}[subsection]{Definition}
\newtheorem{example}[subsection]{Example}
\newtheorem{exercise}[subsection]{Exercise}
\newtheorem{situation}[subsection]{Situation}

\theoremstyle{remark}
\newtheorem{remark}[subsection]{Remark}
\newtheorem{remarks}[subsection]{Remarks}

\numberwithin{equation}{subsection}

% Macros
%
\def\lim{\mathop{\mathrm{lim}}\nolimits}
\def\colim{\mathop{\mathrm{colim}}\nolimits}
\def\Spec{\mathop{\mathrm{Spec}}}
\def\Hom{\mathop{\mathrm{Hom}}\nolimits}
\def\Ext{\mathop{\mathrm{Ext}}\nolimits}
\def\SheafHom{\mathop{\mathcal{H}\!\mathit{om}}\nolimits}
\def\SheafExt{\mathop{\mathcal{E}\!\mathit{xt}}\nolimits}
\def\Sch{\mathit{Sch}}
\def\Mor{\mathop{\mathrm{Mor}}\nolimits}
\def\Ob{\mathop{\mathrm{Ob}}\nolimits}
\def\Sh{\mathop{\mathit{Sh}}\nolimits}
\def\NL{\mathop{N\!L}\nolimits}
\def\CH{\mathop{\mathrm{CH}}\nolimits}
\def\proetale{{pro\text{-}\acute{e}tale}}
\def\etale{{\acute{e}tale}}
\def\QCoh{\mathit{QCoh}}
\def\Ker{\mathop{\mathrm{Ker}}}
\def\Im{\mathop{\mathrm{Im}}}
\def\Coker{\mathop{\mathrm{Coker}}}
\def\Coim{\mathop{\mathrm{Coim}}}

% Boxtimes
%
\DeclareMathSymbol{\boxtimes}{\mathbin}{AMSa}{"02}

%
% Macros for moduli stacks/spaces
%
\def\QCohstack{\mathcal{QC}\!\mathit{oh}}
\def\Cohstack{\mathcal{C}\!\mathit{oh}}
\def\Spacesstack{\mathcal{S}\!\mathit{paces}}
\def\Quotfunctor{\mathrm{Quot}}
\def\Hilbfunctor{\mathrm{Hilb}}
\def\Curvesstack{\mathcal{C}\!\mathit{urves}}
\def\Polarizedstack{\mathcal{P}\!\mathit{olarized}}
\def\Complexesstack{\mathcal{C}\!\mathit{omplexes}}
% \Pic is the operator that assigns to X its picard group, usage \Pic(X)
% \Picardstack_{X/B} denotes the Picard stack of X over B
% \Picardfunctor_{X/B} denotes the Picard functor of X over B
\def\Pic{\mathop{\mathrm{Pic}}\nolimits}
\def\Picardstack{\mathcal{P}\!\mathit{ic}}
\def\Picardfunctor{\mathrm{Pic}}
\def\Deformationcategory{\mathcal{D}\!\mathit{ef}}
\usepackage{cleveref}
\crefname{originallemma}{original lemma}{original lemmas}
\Crefname{originallemma}{Original lemma}{Original lemmas}
\makeatletter
\newlabel{lemma-characterize-constructible-Noetherian}{{lemma-characterize-constructible-Noetherian}{source}{}{sourcepointer-1}{}}
\newlabel{lemma-characterize-constructible-Noetherian@cref}{{[originallemma][0][]lemma-characterize-constructible-Noetherian}{[1][0][]source}{}{}{}}
\makeatother
\title{Constructible subsets of a Noetherian space are open in the constructible topology}
\author{The Stacks Project Authors}
\date{stacks-project:topology.tex:local-snapshot-20261012; GNU FDL 1.2 or later}
\begin{document}
\maketitle
\noindent\textit{Editorial note.} All delivered mathematical text is the author's own, copied from the pinned Stacks Project chapter source sources/stacks-topology/pinned-original-source.tex; no statement or proof step was rewritten, reordered or removed.\par\medskip
\noindent\textit{Editorial note.} Preamble provenance: the Stacks Project keeps one single shared preamble for every chapter; this card takes it directly from the staged sources/stacks-topology/preamble.tex as three verbatim spans (lines 8-12, 14-18 and 154-217), so no derived file was written and nothing was added or redefined. The wrapper class and the base packages are supplied by the assembler.\par\medskip
\noindent\textit{Editorial note.} Bibliography provenance: neither the counted statement nor its proof contains a citation (this lane rescans the whole span before writing the card), so the card carries no bibliography and no entry was reconstructed.\par\medskip
\noindent\textit{Editorial note.} Selection provenance: this card was opened under the file-based control contract (CONTROL.md), with the mechanical pre-card gate run on its statement and proof interval before the card was written.\par\medskip
\noindent\textit{Allowed standard prerequisites (editorial).}\begin{itemize}
\item The characterization of constructible subsets as finite unions of locally closed subsets. Invoked inside the counted proof (source lines 3177-3185). Writes the subset as a finite union of pieces whose openness in the constructible topology is easier to see.
\end{itemize}\medskip
\ifdefined\hypertarget\hypertarget{sourcepointer-1}{}\fi
\noindent\textit{Source reference.} Original lemma lemma-characterize-constructible-Noetherian: sources/stacks-topology/pinned-original-source.tex, source lines 3113--3113. Navigation pointer only: the counted text refers to the Stacks item lemma-characterize-constructible-Noetherian, which is an existing labeled result of this same chapter and is not part of this card. The pointer carries the target's own label as its number and its definition line as its span; its content is not reproduced here.\par\medskip
\begin{lemma}
\label{lemma-constructible-neighbourhood-Noetherian}
Let $X$ be a Noetherian topological space.
Let $x \in X$.
Let $E \subset X$ be constructible in $X$.
The following are equivalent:
\begin{enumerate}
\item $E$ is a neighbourhood of $x$, and
\item for every irreducible closed subset $Y$ of $X$ which contains
$x$ the intersection $E \cap Y$ is dense in $Y$.
\end{enumerate}
\end{lemma}

\begin{proof}
It is clear that (1) implies (2). Assume (2).
Consider the set $\mathcal{S}$ of closed subsets $Y$ of $X$ containing $x$
such that $E \cap Y$ is not a neighbourhood of $x$ in $Y$.
If $\mathcal{S} \not = \emptyset$, then it has a minimal element $Y$
as $X$ is Noetherian. Suppose $Y = Y_1 \cup Y_2$ with two smaller nonempty
closed subsets $Y_1$, $Y_2$. If $x \in Y_i$ for $i = 1, 2$, then $Y_i \cap E$
is a neighbourhood of $x$ in $Y_i$ and we conclude $Y \cap E$ is a
neighbourhood of $x$ in $Y$ which is a contradiction. If $x \in Y_1$ but
$x \not\in Y_2$ (say), then $Y_1 \cap E$ is a neighbourhood of $x$ in
$Y_1$ and hence also in $Y$, which is a contradiction as well.
We conclude that $Y$ is irreducible closed. By assumption (2) we see that
$E \cap Y$ is dense in $Y$. Thus $E \cap Y$ contains an open $V$ of $Y$, see
Lemma \ref{lemma-characterize-constructible-Noetherian}.
If $x \in V$ then $E \cap Y$ is a neighbourhood of $x$ in $Y$ which
is a contradiction. If $x \not \in V$, then $Y' = Y \setminus V$ is a
proper closed subset of $Y$ containing $x$. By minimality of $Y$
we see that $E \cap Y'$ contains an open neighbourhood $V' \subset Y'$
of $x$ in $Y'$. But then $V' \cup V$ is an open neighbourhood of $x$
in $Y$ contained in $E$, a contradiction.
This contradiction finishes the proof of the lemma.
\end{proof}

\end{document}
